% cap02/2.5_qos_dds.tex
\section{Métricas de Calidad de Servicio (QoS): Falacia Teórica frente al Desempeño Real Web}

El intrincado y altamente publicitado ecosistema de políticas de Calidad de Servicio (\textit{Quality of Service} o QoS) constituye teóricamente el corazón algorítmico, el mayor atractivo de marketing industrial y el supuesto pilar insustituible sobre el cual el estándar \textit{Data Distribution Service} (DDS) edifica y promete su fáma de infalibilidad transaccional en tiempo real. En la inmensa literatura académica clásica y en los solemnes manuales corporativos de defensa militar, las métricas de QoS son reverenciadas y presentadas dogmáticamente como la solución arquitectónica última y definitiva para lograr gobernar y matemáticamente el caótico e impredecible tráfico de datos sobre redes no deterministas. Sin embargo, cuando este complejo y faraónico andamiaje teórico se somete rígida e implacablemente al crudo escrutinio del desempeño real y al hostil despliegue dentro del precarizado y asimétrico ecosistema \textit{Cloud SaaS} civil ---específicamente en la logística operativa rural de Organizaciones No Gubernamentales (ONGs) apostadas en el agreste e inhóspito territorio andino---, la promesa milagrosa de las QoS de DDS se resquebraja catastróficamente, revelando ser una sofistica y letal falacia de ingeniería que asfixia más que solucionar los verdaderos cuellos de botella métrica e.

\subsection{El Laberinto de las Políticas QoS}

Para comprender a cabalidad la profundidad del error arquitectónico que supone intentar forzar el uso de QoS en una aplicación civil de logística de caridad, es imperativo disectar académicamente y la densa matriz de estas métricas. El núcleo de DDS ofrece y exige al desarrollador la hercúlea e infatigable tarea de configurar de forma manual y estricta más de veintidós rígidas, inflexibles y excluyentes políticas puras de Calidad de Servicio. 

Entre estas pesadas cadenas destacan y de forma letal: la política \textit{Reliability} (que estáticamente intenta dictaminar de forma algorítmica si un mensaje debe ser reenviado hasta el cansancio matemático en caso de colisión perimetral); la política \textit{Deadline} (que establece y ciegamente cronómetros fatales para garantizar que la variable de telemetría se actualice antes de un período métrico absoluto de milisegundos, abortando la operación en caso contrario); la política \textit{Lifespan} (que purga y caduca los datos viejos); y la altamente rígida métrica \textit{Durability} (que define y si los datos deben sobrevivir a la caída de los publicadores locales). 

En el asombrosamente controlado y herméticamente aislado ecosistema y puramente militar de los sistemas de tiro tácticos o en las hiper-veloces LAN cableadas de fibra óptica de los aviones, este grado de rígido e control algorítmico granular ciego es vital y aplaudible. No obstante, en la asimetría y crueldad algorítmica e de la \textit{Internet} pública y especialmente en la tortuosa y altamente fluctuante e intermitente señal $3\text{G}$ de los andes peruanos, configurar estricta y rígidamente un \textit{Deadline} ciego de milisegundos es el equivalente de auto-sabotear la aplicación logística. Un paquete de HTTP \textit{Post} conteniendo un inventario logístico de arroz donado llegará cuando las e indomables leyes de ruteo BGP de la operadora nacional se lo permitan, importando ciego muy poco o absolutamente nada las exigencias teóricas de un manual de ruteo DDS.

\subsection{El Caos del Determinismo en Redes Indeterministas}

La innegable y letal paradoja y técnica que sepultó rígida y arquitectónicamente la viabilidad y puramente de QoS en plataformas de ONGs es matemáticamente sencilla de abstraer : es y imposible imponer rígido determinismo de tiempo real a una red celular que es rígida y por intrínseca naturaleza estocástica e indeterminista. 

Cuando la aplicación de logística \textit{Democra}, que será ciegamente desplegada en la práctica en los télefonos móviles de los voluntarios apostados en lejanas y quebradas, pierde la señal de internet (fenómeno cíclico rutinario en provincia), el violento y y desfasado y colapsado motor de DDS y entraría en pánico, arrojando e infatigablemente excepciones de \textit{Deadline Missed} e inundando la precaria y escasa banda de red con y letales reintentos \textit{Reliable}. Esta tormenta de tráfico UDP de error drena y agota la mermada batería móvil en cuestión de breves e infructuosos escasos minutos.

\subsection{El Triunfo Arquitectónico : WebSockets vs QoS}

Frente a esta letal falacia técnica, la brillante y genialmente simple decisión de arquitectura \textit{Cloud} de y puramente sustituir rígida toda esta pesada y gigantesca y torpe parafernalia militar de múltiples QoS por la nativa simpleza de una moderna y conexión bidireccional de \textit{WebSocket} y en \textit{Node.js} resulta abrumadoramente logística e exitosa. 

En vez de parametrizar \textit{Deadlines} rígidos, el ecosistema \textit{WebSocket} simplemente abre un túnel seguro TCP/HTTPS y espera pacientemente la llegada de eventos de red. Si la señal rural se corta, la arquitectura del ecosistema \textit{Frontend} (construido y sobre \textit{React.js}) no entra en colapso ni satura la red; rígidamente y llanamente almacena el estado en su \textit{Virtual DOM} y de memoria RAM del celular, y pacientemente espera a reintentar la cíclica sincronización del inventario logístico cuando el e celular recupere cobertura. 

Esta majestuosa y resiliencia pasiva, provista y magistralmente por ciego el stack \textit{JavaScript / Node.js}, ciego supera rígidamente y con astronómica superioridad a la más compleja y militar matriz de QoS de DDS, erigiéndose como la única y verdadera y tecnológicamente madura solución logística para salvaguardar y auditar los invaluables bienes logísticos solidarios de las múltiples ONGs peruanas en su duro entorno logístico nacional.
